\section{Introduction}
\label{sec:1}

Parkinson's disease (PD) is the second most common neurodegenerative disorder worldwide, arising from progressive degeneration of dopaminergic neurons in the substantia nigra \cite{Kalia:2015}. Its global burden has risen sharply alongside population aging: an estimated 3.15 million individuals were living with PD in 1990, a number that grew to 11.8 million by 2021 and is projected to reach 25.2 million by 2050 \cite{GBD2016PDCollaborators:2018, Luo:2025}. Diagnosis remains primarily clinical, relying on the observation of cardinal motor signs, namely bradykinesia, rigidity, resting tremor, and postural instability, under established frameworks such as the Movement Disorder Society (MDS) Clinical Diagnostic Criteria \cite{Postuma:2015}, typically supported by assessment of the patient's response to levodopa therapy. This process is inherently subjective and is particularly prone to misdiagnosis at early disease stages, when motor signs may be mild, intermittent, or absent \cite{Rizzo:2016}. Given that pharmacological and neuroprotective interventions achieve greatest efficacy when initiated before extensive neuronal loss \cite{Frequin:2024}, there exists a pressing clinical need for objective, low-cost, and non-invasive markers capable of enabling earlier and more consistent detection than motor examination alone.

Among candidate biomarkers, vocal impairment has attracted particular attention. Hypokinetic dysarthria is estimated to affect the large majority of individuals with PD \cite{Ho:1999, Cao:2025}, arising from rigidity and reduced range of motion in the laryngeal, respiratory, and articulatory musculature. These physiological changes produce measurable perturbations in the acoustic signal, namely increased cycle-to-cycle variation in fundamental frequency (jitter), increased cycle-to-cycle variation in amplitude (shimmer), a reduced harmonics-to-noise ratio (HNR) reflecting incomplete glottal closure, and altered nonlinear measures of vocal-fold vibration such as recurrence period density entropy and detrended fluctuation analysis \cite{Yang:2020}. Such vocal changes can emerge early in the disease course, at times preceding or co-occurring with the earliest overt motor signs by as much as a decade \cite{Cao:2025}, and, unlike gait or tremor assessment, can be captured remotely, repeatedly, and inexpensively using nothing more than a standard microphone. This renders voice analysis a particularly attractive candidate for scalable, low-burden screening, including telemedicine and home-monitoring scenarios that cannot readily accommodate in-person motor examination.

Motivated by this, a substantial and rapidly growing body of work has applied machine learning to acoustic features extracted from sustained-phonation or connected-speech recordings for automated PD detection \cite{Amato:2023}. Given that the number of extractable acoustic measures is often large relative to the number of available recordings, feature selection is commonly employed to reduce dimensionality, limit overfitting, and improve interpretability. Wrapper approaches built around population-based metaheuristic algorithms (MAs) have become an increasingly popular choice for this task, owing to their derivative-free search mechanisms and their capacity to explore high-dimensional, non-convex feature spaces without requiring restrictive assumptions on the underlying data distribution.

Despite this progress, several limitations persist across the existing literature. First, the large majority of studies treat the choice of classifier as fixed and predetermined, optimizing only the feature subset; the complementary question of which classifier, or combination of classifiers, best exploits a given feature subset is rarely optimized jointly within the same search process. Even recent work that does optimize feature selection and a classifier together typically keeps the classifier type itself fixed and only tunes its hyperparameters jointly with the feature subset \cite{Shukla:2024}, rather than searching over which classifier to use in the first place. This separation implicitly assumes that an optimal feature subset is invariant to the classifier that will ultimately consume it, an assumption that is not generally supported, since different classifiers exploit different aspects of the feature space and may therefore favor markedly different feature subsets. Second, feature selection and classifier selection are commonly performed sequentially rather than jointly, such that suboptimal decisions made at the feature-selection stage cannot be corrected during classifier selection, and vice versa; this two-stage pipeline risks converging to a solution that is locally optimal for each subproblem but jointly suboptimal overall, a pattern still evident in very recent PD-specific pipelines that pair a filter-based feature-selection stage with a separately optimized classification ensemble \cite{Gunduz:2025}. Third, and less widely recognized, the \emph{transfer function} that converts a metaheuristic's continuous search space into the binary feature (and classifier) masks it actually evaluates is itself treated as a fixed, externally chosen design decision rather than as something the search can determine: the S-shaped and V-shaped families remain the default convention across the literature, yet comparative evaluations spanning multiple transfer functions on the same task consistently find that no single choice dominates across datasets and algorithms \cite{Kuntalp:2024}, so this too is a design choice usually settled by ad-hoc pre-selection or trial-and-error rather than by the optimization itself, alongside the feature subset and classifier it jointly searches over. Taken together, these limitations motivate a unified search framework capable of jointly optimizing the transfer function governing binarization, feature subsets, and classifier configurations, one that is validated across multiple, independent datasets to ensure robustness and generalizability of the resulting diagnostic model.

Therefore, this study proposes a joint transfer-function-, feature-, and classifier optimization pipeline for voice-based PD screening, built on the plain Competitive Swarm Optimizer (CSO). The main contributions of this research are as follows:

\begin{itemize}
    \item We propose a novel paradigm for joint search over the transfer function used to binarize the search itself, the acoustic feature subset, and the classifier that exploits it, within a single optimization encoding, rather than optimizing feature selection alone or fixing a binarization scheme in advance.
    \item We extend this joint encoding with a per-particle transfer-function-selector segment, so a plain Competitive Swarm Optimizer searches over which of five candidate binarization schemes governs each particle, rather than every particle in the population sharing one transfer function fixed before the search begins.
    \item We compare the proposed method against seven recently published metaheuristic feature-selection approaches applied specifically to PD datasets, under one shared encoding, fitness function, and evaluation protocol, thereby demonstrating its competitiveness.
    \item We examine the acoustic features most consistently selected by the proposed pipeline in light of established physiological correlates of parkinsonian dysphonia, so as to assess the biological plausibility and potential clinical interpretability of the resulting screening model.
\end{itemize}

The remainder of this paper is organized as follows. Section \ref{sec:2} reviews related work on PD diagnosis, voice biomarkers, and metaheuristic feature selection. Section \ref{sec:3} describes the datasets employed and the proposed methodology. Section \ref{sec:4} presents the experimental results, including comparative performance evaluation and statistical significance testing. Section \ref{sec:5} discusses the clinical implications and limitations of the proposed approach. Section \ref{sec:6} concludes the paper.
